\section{Executive summary and exact scope}

\subsection{The problem exposed by the repository}

The repository separates two very different achievements.  The directory
\code{ConciseIPRA} contains a Lean formalization of ordinary first-order
intuitionistic primitive recursive arithmetic and its bi-interpretation with a
small native scalar calculus.  In particular, the repository provides the
object theory \(\PRAAx\), derivations \(\Prf\), quantifier-free induction
\code{PRA.qfInd}, and the transfer theorem \code{pra\_iff\_native}.  By contrast,
the theorem packets under \code{docs/proofs} are research reports: most contain
mathematical arguments, executable finite checks, and Lean plans, but none is a
completed derivation of its advertised theorem in \(\Prf\,\PRAAx\).

The recurring difficulty is therefore not merely ``how to code real numbers''
or ``how to avoid excluded middle.''  It is how to organize a proof so that
all of the following are simultaneously explicit:

\begin{enumerate}[label=(\arabic*)]
\item the finite information a construction reads;
\item the primitive-recursive algorithm or certificate checker operating on
      that information;
\item a quantifier-free invariant suitable for \(\ipra\)'s induction;
\item the modulus, separation margin, height bound, finite net, or covering
      data that makes an apparently infinitary step finite;
\item the exact object-language formula and derivation to be transported to
      the native calculus;
\item the boundary between a completed verifier and an unconstructed producer.
\end{enumerate}

The existing foundation packets already contain many of these ingredients.
The present report unifies them and adds two missing interfaces: a syntactic
proof-producing compiler and a relational refinement layer for non-unique
approximate witnesses.

\subsection{The proposed architecture}

The framework has five layers.

\paragraph{Layer 1: fixed numerical presentations.}
Every object family is represented by natural-number parameters and fixed
primitive-recursive evaluators.  A represented object is not an arbitrary
program supplied at run time.  If an external computation is not represented
by a fixed evaluator, its finite trace is supplied as a certificate and checked
by a fixed primitive-recursive verifier.  This distinction avoids the false
idea of a total universal primitive-recursive evaluator.

\paragraph{Layer 2: Boolean finite contracts.}
All bounded quantifiers, list checks, finite sums, maxima, residual tests, and
search predicates are compiled to primitive-recursive Boolean functions.
Global name hypotheses remain ordinary assumptions in the context; induction
uses only equations saying that a Boolean accumulator equals one.  Precision is
kept free during the finite proof and universally quantified only afterward.

\paragraph{Layer 3: proof-producing traces.}
A bounded computation is presented by fixed functions
\(\mathsf{init},\mathsf{step},\mathsf{bound},\mathsf{inv},\mathsf{accept}\).
Given object-language derivations of the base, local step, and final soundness
formulas, a generic theorem constructs an object-language derivation of the
whole run.  The construction uses \code{PRA.qfInd}, the defining equations of
primitive recursion, equality substitution, and ordinary natural deduction.
Thus the result is not merely an ambient theorem about a Lean function.

\paragraph{Layer 4: finitary relational presentations.}
A challenge may admit several correct approximate answers.  The framework
therefore permits a finite relation \(\Good(p,k,a)\) rather than insisting on a
single canonical \(a\).  A compatibility relation connects answers at adjacent
precision levels.  Bounded backward dynamic programming determines which
answers can be extended to a finite horizon and constructs a compatible finite
chain whenever one exists.  This is enough for all fixed-precision theorems and
avoids a hidden choice principle.

\paragraph{Layer 5: determinization only from quantitative data.}
When a downstream theorem genuinely needs one coherent name, it must consume a
quantitative determinacy hypothesis.  The shrinking-viability theorem assumes
that all candidates viable to a stated horizon lie in a small observational
ball.  The least viable candidate then defines a primitive-recursive name, and
the shrinking condition proves regularity.  Root separation, strong convexity,
isolating intervals, contraction estimates, or projective gaps are instances
of this hypothesis.  Qualitative uniqueness alone is not silently converted
into a modulus.

\subsection{Principal theorems}

The main general results proved in this report are:

\begin{enumerate}[label=\textbf{T\arabic*.},leftmargin=4em]
\item \textbf{Bounded Booleanization.} Every fixed bounded formula over
      primitive-recursive terms has a Boolean evaluator whose correctness is
      provable in \(\ipra\).
\item \textbf{Finite-trace compiler.} Embedded base, step, and final
      derivations yield an embedded theorem about the complete bounded run.
\item \textbf{Contract composition.} Proof-producing transducers compose,
      including adaptive finite query trees.
\item \textbf{Finite witness amalgamation.} A bounded backward viability
      computation constructs a compatible chain through finite candidate
      levels whenever such a chain exists.
\item \textbf{Shrinking viability.} Finite candidate relations plus a
      quantitative shrinking horizon define a coherent primitive-recursive
      approximate name.
\item \textbf{Guarded extensionality.} Operations with precision transformers
      descend to observational setoids, while partial operations are totalized
      and guarded by explicit domain or apartness certificates.
\item \textbf{Finite-net transfer.} A modulus and a finite net turn local
      rational checks into uniform conclusions on a represented compact set.
\item \textbf{Height bound/list equivalence.} For a decidable coded predicate,
      an explicit height bound and a complete finite list are constructively
      interconvertible; a mere cardinality bound is not enough.
\end{enumerate}

\subsection{What is and is not resolved}

The framework resolves several obstacles at the level at which they are
actually logical or architectural.  It gives a reusable path from finite
algorithms to \(\Prf\,\PRAAx\), a correct treatment of all-precision hypotheses,
a generic finite-search library, and a non-choice-based interface for
non-unique approximations.  It also gives standard positive replacements for
exact sign, root, axis, pivot, and branch selection.

It does not manufacture missing mathematics.  In particular, it does not
produce the unknown or ineffective height bounds needed for general Faltings
or Siegel statements, rational covering representatives for arbitrary
Mordell--Weil descent, or all constants in the strongest prime-distribution
claims.  Instead, those results are factored into a uniform finite verifier and
an explicit producer obligation.  This is progress: the verified downstream
arithmetic can be completed once a producer is supplied, and no search cutoff
is smuggled in through a qualitative theorem.

The included Lean files are detailed blueprints against the current repository
API.  They contain theorem statements, formula constructors, proof dependency
order, and exact points where \code{PRA.qfInd}, \code{Prf.subst},
\code{Prf.trans}, and \code{pra\_iff\_native} are used.  Lean was not available
in the artifact environment, so this archive does not claim that those files
have been elaborated or kernel-checked.
